\cotefile{157-2}
\scanpage{99}
\mspage{43}
% archive sheet number: 98
les inclusions

\[
 X_s\longrightarrow {}_{s\backslash}X_{\deleted{\unsure{/s|s}}},\quad X_s\longrightarrow X_{/s}\quad\deleted{\unsure{{}_{s\backslash}X_s}}
\]

et de $=$ pour $Y$, sont des $W_{\infty}$-équiv. donc des $W$-équiv.
% slanted
\margin{comparer déf.~3 p.~117 \unsure{du|de} (I)}
(Il suffirait $=$ que $X,Y$ soient des $W$-fibrations \deleted{\unsure{faibles}} sur $S$*, i.e. \deleted{par déf.} que les foncteurs précédents soient des $W$-équivalences, et que \deleted{ils soient des \unsure{équiv.}} \inserted{la $=$ condition soit satisfaite} après tt chgt. de base sur $S$), l'argument \unsure{dans le cas} lisse resp. propre montre déjà que

\diagram{\begin{tikzcd}[ampersand replacement=\&,column sep=small,row sep=small]
 c \& b\arrow[l,Rightarrow]\arrow[r,Leftrightarrow]\arrow[d,Rightarrow] \& c'\\
 \& a \&
\end{tikzcd}}

% slanted
\margin{valable en supposant seulement $X,Y$ des $W$-\unsure{fibrissimes|fibrissures}.}

, et que la $W$-équivalence par fibres \deleted{implique} (ou la $W$-\deleted{\unsure{équiv.}} locale sur $S$, resp. bilocale sur $(\,)$) implique déjà la $W$-équivalence \uline{universelle} sur $S$.

La question $a)\Rightarrow b)$ est par contre \unsure{nettement|nullement} plus délicate, et on a vu que ce n'est \deleted{pas} nullement le cas pour tout \unsure{localisateur|localiseur} fondamental -- sans doute pas pour celui associé au dérivateur $\mathrm{Der}(\mathbb Z\text{-}\unsure{\mathrm{Mod}|\mathrm{Modules}})$, en tous cas pas à celui associé à $\mathrm{Der}(k\text{-}\mathrm{Mod})$, quand $2$ est inv. dans $k$. (Il suffirait $=$ \deleted{\unsure{que un}} que $\exists n\in\mathbb N,\ n\geq 2$, tel que $n$ inv. dans $k$ -- le contrexemple typique est

\diagram{\begin{tikzcd}[ampersand replacement=\&,row sep=tiny,column sep=small]
 S^1\arrow[r,"n"]\arrow[d,equal] \& S^1\arrow[d,equal]\\
 X \& Y=S
\end{tikzcd}}

, le rev. étale d'ordre $n$ d'un cercle. \inserted{L'implication $a\Rightarrow b$} \deleted{\unsure{C'}}est \inserted{donc} une condition très forte sur $W$, je vais la noter $\mathrm L\,5$ (l'ancienne L5 \unsure{``$3^{\mathrm e}$''}, cf page 2, \deleted{\unsure{que}} ne \unsure{méritant|méritait} pas de survivre, car c'est \unsure{une} côte mal taillée)

% slanted
\margin{Dans \uline{XII} (p.~12) on voit que si $W$ satisfait à L1, \inserted{L3 quater}, L4 bis (cf. p.~\unsure{suivante}), L5, alors \unsure{vraiment} $\boxed{W=W_{\infty}}$ ! Donc c'est vraiment en \unsure{supprimant|supposant} \unsure{l'hyp.} \unsure{L5} qu'on \unsure{s'aventure} \unsure{hors} du train-train}

\uline{L 5} Soit

\diagram{\begin{tikzcd}[ampersand replacement=\&,column sep=small,row sep=tiny]
 X\arrow[rr,"f"]\arrow[dr] \& \&Y\arrow[dl]\\
 \&S\&
\end{tikzcd}}

dans $\mathrm{Cat}$, avec $X,Y$ des \unsure{Hot}-fibrés sur $S$. Si $f$ est une \deleted{\unsure{Hot}} $W$-équiv. \inserted{(``grossière'' p.r. à $S$)} alors c'est une $W$-équivalence par fibres.

On a vu qu'à un $\varepsilon$ près (\unsure{``lemme''} p.~23), L5 est une conséquence de

\uline{L 6} \deleted{\unsure{le foncteur}} Soit $f:X\longrightarrow Y$ une $W$-équivalence, $Y'\longrightarrow Y$ \deleted{une} \inserted{morphisme parfait, (\unsure{c.à.d.} lisse et propre)} \deleted{\unsure{Hot-fibration}}, alors $f':X'=X\times_Y Y'\longrightarrow Y'$ est une $W$-équiv.

